% CVPR 2026-style manuscript (self-contained draft)
\PassOptionsToPackage{table}{xcolor}
\documentclass[10pt,twocolumn,letterpaper]{article}

% \usepackage{cvpr}
\usepackage[review]{cvpr}
% \usepackage[pagenumbers]{cvpr}

\usepackage{amsmath,amssymb}
\usepackage{graphicx}
\usepackage{booktabs}
\usepackage{multirow}
\usepackage{array}
\usepackage{tabularx}
\usepackage{threeparttable}
\usepackage{subcaption}
\usepackage{algorithm}
\usepackage{algorithmic}
\usepackage{enumitem}
\usepackage{url}

\definecolor{cvprblue}{rgb}{0.21,0.49,0.74}
\usepackage[pagebackref,breaklinks,colorlinks,citecolor=cvprblue]{hyperref}

\newcommand{\bbR}{\mathbb{R}}
\newcommand{\method}{\textsc{BASS}}

\def\paperID{9}
\def\confName{CVPR}
\def\confYear{2026}

\title{Budget-Aware Scene Subset Selection for Efficient NeRF Benchmarking at Nutrition5k Scale}

\author{Ahmad AlMughrabi\\
Universitat de Barcelona, Spain\\
{\tt\small ahmad.almughrabi@ub.edu}
\and
Flavi\`a Serret\\
Universitat de Barcelona, Spain\\
{\tt\small flserrem7@alumnes.ub.edu}
\and
Ricardo Marques$^*$\\
Grup de Tecnologies Interactives (GTI),\\
Universitat Pompeu Fabra (UPF), Spain\\
{\tt\small ricardo.marques@upf.edu}
\and
Petia Radeva$^*$\\
Universitat de Barcelona, Spain\\
Institut de Neuroci\`encies, Barcelona\\
{\tt\small petia.ivanova@ub.edu}
}

\begin{document}
\maketitle

\def\thefootnote{*}\footnotetext{Equal supervision.}
\def\thefootnote{\arabic{footnote}}

\begin{abstract}
Benchmarking neural radiance field (NeRF) methods at Nutrition5k scale (\(\sim\)5,000 scenes) is computationally expensive because each method is trained and evaluated per scene. We address budget-aware scene subset selection: choosing the smallest subset that preserves benchmark conclusions under compute constraints. We present \method, a practical pipeline combining feature-space representativeness, difficulty-aware coverage, and explicit metric-matching validation against full-dataset PSNR/SSIM statistics. The system includes risk-aware budget auto-selection using a lower confidence bound (LCB) of pass probability, optional per-cluster gap constraints, joint \(k\times\)budget search, and holdout selection/tune/test protocol. On current Nutrition5k pipeline outputs (3,522 scenes, 57 features), \method produces compact subsets that satisfy strict global matching thresholds (\(|\Delta\text{PSNR}|\le 0.5\), \(|\Delta\text{SSIM}|\le 0.01\)) while reducing benchmark size substantially. We release a reproducible implementation with end-to-end report generation.
\end{abstract}

\section{Introduction}
Neural radiance fields (NeRFs) deliver high-quality novel view synthesis but remain expensive for large benchmark suites because optimization is typically scene-specific \cite{mildenhall2020nerf,barron2021mipnerf,muller2022instantngp}. At Nutrition5k scale, exhaustive evaluation across many methods can be prohibitive in GPU-hours.

The key practical question is: \emph{what is the minimum scene subset that preserves benchmark conclusions?} A useful subset must preserve distributional coverage and method-comparison behavior. Naive random subsets can miss failure regimes; pure representativeness can still under-cover difficult strata.

We build on an implemented scene-feature pipeline and cast subset selection as a budget-aware optimization with explicit post-selection validation. The resulting system, \method, supports two operating modes: (i) fixed clustering granularity with budget sweep; and (ii) joint \(k\times\)budget search. In both modes, recommendations are made with risk-aware rules based on pass-rate confidence bounds, not only mean performance.

\paragraph{Contributions.}
\begin{enumerate}[leftmargin=*,itemsep=0pt]
\item A practical budget-aware NeRF subset selection framework with explicit global and per-cluster metric-matching constraints.
\item A risk-aware auto-selection strategy using Wilson lower confidence bounds (LCB) on joint pass rate.
\item Joint \(k\times\)budget recommendation, holdout selection/tune/test protocol, and metric-aware subset refinement by cluster-preserving swaps.
\item Reproducible tooling for extraction, sweeps, validation, visualization, and final report packaging.
\end{enumerate}

\section{Related Work}
\paragraph{NeRF and acceleration.}
The original NeRF formulation \cite{mildenhall2020nerf} enabled photorealistic view synthesis from posed images, followed by quality and efficiency improvements such as mip-NeRF \cite{barron2021mipnerf} and Instant-NGP \cite{muller2022instantngp}. Despite acceleration, cross-scene benchmarking remains costly.

\paragraph{Structure-from-motion signals.}
Scene reconstruction pipelines such as COLMAP \cite{schoenberger2016sfm,schoenberger2016mvs} provide geometric diagnostics (track lengths, reprojection error, sparse point statistics) that are useful proxies for scene difficulty.

\paragraph{Subset selection and submodularity.}
Representative subset selection is classically addressed by facility-location objectives and submodular maximization \cite{lin2011submodular,mirzasoleiman2015lazier}, clustering medoids, or k-center/farthest-first selection \cite{gonzalez1985kcenter}. Under cost constraints, knapsack-style greedy approximations are widely used.

\paragraph{Rank fidelity and uncertainty.}
Benchmark consistency often depends on method ordering. Rank correlations (Spearman, Kendall) and bootstrap uncertainty estimates are standard tools for stability analysis.

\section{Method}
\subsection{Problem Definition}
Let scenes be \(\mathcal{D}=\{1,\ldots,N\}\), each with features \(\mathbf{x}_i\in\bbR^d\) and optional runtime proxy \(c_i\). We seek a subset \(\mathcal{S}\subseteq\mathcal{D}\) that preserves benchmark behavior under size/budget constraints.

For source \(s\in\{\text{OC},\text{FI}\}\) and metric \(m\in\{\text{PSNR},\text{SSIM}\}\), define global gap:
\begin{equation}
\Delta_{s,m}(\mathcal{S}) = \mu_{s,m}(\mathcal{S}) - \mu_{s,m}(\mathcal{D}).
\end{equation}
A subset passes global constraints if
\begin{equation}
|\Delta_{s,\text{PSNR}}|\le\tau_{\text{PSNR}},\quad
|\Delta_{s,\text{SSIM}}|\le\tau_{\text{SSIM}},\;\forall s.
\end{equation}

Optional per-cluster constraints enforce local fidelity:
\begin{equation}
\max_c |\Delta^{(c)}_{s,\text{PSNR}}|\le\tau^{(c)}_{\text{PSNR}},\quad
\max_c |\Delta^{(c)}_{s,\text{SSIM}}|\le\tau^{(c)}_{\text{SSIM}}.
\end{equation}

\subsection{Feature Backbone and Candidate Generation}
Our feature pipeline extracts camera/intrinsics, image quality, mask geometry, texture density, COLMAP sparse geometry, and pose coverage statistics per scene; normalization uses log-transform candidates, robust/standard scaling, and clipping.

For a cluster setting \(k\) and per-cluster budget \(b\), candidate subsets are formed by balanced sampling of \(b\) scenes per cluster. This yields \(|\mathcal{S}|\approx k\cdot b\) (subject to cluster sizes).

\subsection{Pass-Rate Estimation and Risk-Aware Selection}
For each \((k,b)\), we run \(M\) balanced simulations. Let \(I_t\in\{0,1\}\) indicate whether candidate \(t\) passes all constraints. The joint pass rate is
\begin{equation}
\hat p(k,b)=\frac{1}{M}\sum_{t=1}^M I_t.
\end{equation}
To avoid optimistic selection, we use Wilson lower confidence bound:
\begin{equation}
\mathrm{LCB}_{\alpha}(k,b)=\mathrm{WilsonLower}(\hat p(k,b),M,\alpha).
\end{equation}
Auto-selection minimizes subset size while satisfying \(\mathrm{LCB}_{\alpha}\ge p_{\min}\); if no candidate satisfies target, fallback selects highest \(\mathrm{LCB}_{\alpha}\).

\subsection{Best-Manifest Objective}
Among simulated candidates per \((k,b)\), we keep one manifest minimizing normalized mismatch:
\begin{equation}
J(\mathcal{S})=
\sum_{s\in\{\text{OC},\text{FI}\}}
\left(
\frac{|\Delta_{s,\text{PSNR}}|}{\tau_{\text{PSNR}}} +
\frac{|\Delta_{s,\text{SSIM}}|}{\tau_{\text{SSIM}}}
\right)
+\lambda_{\text{rt}}\,\widetilde{C}(\mathcal{S}).
\end{equation}
\(\widetilde{C}\) is an optional normalized runtime term.

\subsection{Joint \texorpdfstring{\(k\times\)}{k x}Budget Search}
Instead of fixing \(k\), we evaluate a grid \(\mathcal{K}\times\mathcal{B}\). Each pair returns \((\hat p,\mathrm{LCB}_{\alpha},|\mathcal{S}|,J,\text{manifest})\). Final recommendation applies the same risk-aware rule over the full grid.

\subsection{Holdout Protocol}
To reduce selection bias:
\begin{enumerate}[leftmargin=*,itemsep=0pt]
\item Build stratified split by cluster: selection/tune/test.
\item Select \((k,b)\) on selection split only.
\item Generate split-specific manifests for tune and test with selected budget profile.
\item Report tune/test pass and metric gaps.
\end{enumerate}

\subsection{Metric-Aware Refinement}
Given an initial manifest, we run local cluster-preserving swaps:
\begin{equation}
(i_{\text{old}},i_{\text{new}})\in\text{same cluster} \quad \text{if } J \text{ decreases}.
\end{equation}
This improves metric matching while preserving cluster counts and subset size.

\begin{algorithm}[t]
\caption{\method: Risk-Aware Budget Selection (per \((k,b)\))}
\label{alg:bass}
\begin{algorithmic}[1]
\REQUIRE Cluster mapping, metrics tables, budget \(b\), simulations \(M\), thresholds.
\FOR{\(t=1,\dots,M\)}
  \STATE Sample balanced subset (\(b\) scenes per cluster)
  \STATE Compute global/per-cluster metric gaps
  \STATE Mark pass/fail under thresholds
\ENDFOR
\STATE Estimate \(\hat p\) and \(\mathrm{LCB}_{\alpha}\)
\STATE Choose best manifest minimizing normalized objective \(J\)
\RETURN statistics and best manifest
\end{algorithmic}
\end{algorithm}

\section{Experimental Setup}
\label{sec:exp}

\subsection{Data and Features}
Current pipeline outputs contain 3,522 valid scenes and 57 features per scene. Features combine capture, quality, mask, texture, sparse geometry, and pose-coverage descriptors.

\subsection{Thresholds and Validation}
Unless stated otherwise, we use:
\begin{equation}
\tau_{\text{PSNR}}=0.5,\qquad \tau_{\text{SSIM}}=0.01.
\end{equation}
We report global OC/FI PSNR/SSIM gaps and optional per-cluster maximum gap checks.

\subsection{Selection Configurations}
We evaluate:
\begin{itemize}[leftmargin=*,nosep]
\item fixed-\(k\) budget sweep (e.g., \(k=6\), budgets \(b\in\{2,4,6,8\}\)),
\item joint \(k\times b\) sweep over previously computed \(k\)-runs,
\item holdout protocol with selection/tune/test split,
\item optional refinement via local swaps.
\end{itemize}

\subsection{Implementation and Reproducibility}
All stages are implemented as scripts: feature extraction/reporting, clustering analysis, budget sweep, joint sweep, validation, refinement, holdout protocol, and final report pack generation. Artifacts include CSV/JSON/Markdown summaries and interactive HTML visuals.

\section{Results}

\subsection{Fixed-\texorpdfstring{\(k\)}{k} Budget Sweep and Auto-Selection}
For fixed \(k=6\), we sweep budgets \(b\in\{2,4,6,8\}\) with 400 balanced simulations. Risk-aware auto-selection (LCB criterion, max subset size 48) recommends \(b=8\) (48 scenes), with estimated pass rate 0.1075 and LCB 0.0808.

\begin{table}[t]
\centering
\caption{Fixed-\(k=6\) budget sweep summary (global thresholds: \(|\Delta\mathrm{PSNR}|\le 0.5\), \(|\Delta\mathrm{SSIM}|\le 0.01\)).}
\label{tab:budget_sweep}
\begin{tabular}{rcccc}
\toprule
\(b\) & \(|\mathcal S|\) & pass rate & LCB & recommended \\
\midrule
2 & 12 & 0.0075 & 0.0026 & No \\
4 & 24 & 0.0350 & 0.0210 & No \\
6 & 36 & 0.0600 & 0.0406 & No \\
8 & 48 & 0.1075 & 0.0808 & Yes \\
\bottomrule
\end{tabular}
\end{table}

\subsection{Global Metric-Matching Quality}
The recommended 48-scene subset passes all global and per-cluster checks in full-data validation.

\begin{table}[t]
\centering
\caption{Global metric gaps for recommended 48-scene subset.}
\label{tab:global_gaps}
\begin{tabular}{lcccc}
\toprule
Source & Metric & Full mean & Subset mean & \(|\Delta|\) \\
\midrule
FI & PSNR & 16.5230 & 16.6118 & 0.0888 \\
FI & SSIM & 0.7433 & 0.7454 & 0.0021 \\
OC & PSNR & 22.8016 & 22.6945 & 0.1071 \\
OC & SSIM & 0.7848 & 0.7878 & 0.0029 \\
\bottomrule
\end{tabular}
\end{table}

\subsection{Joint \texorpdfstring{\(k\times\)}{k x}Budget Selection}
Joint search over \((k,b)\) candidates (example run) selects \(k=24\), \(b=2\), total subset size 46, with pass rate 0.1125 and LCB 0.0851, satisfying target LCB 0.08.

\begin{table}[t]
\centering
\caption{Top joint \(k\times b\) candidate from recommendation artifact.}
\label{tab:joint_reco}
\begin{tabular}{ccccc}
\toprule
\(k\) & \(b\) & \(|\mathcal S|\) & pass rate & LCB \\
\midrule
24 & 2 & 46 & 0.1125 & 0.0851 \\
\bottomrule
\end{tabular}
\end{table}

\subsection{Holdout Protocol}
With stratified selection/tune/test protocol, the selection split recommends \(k=12, b=4\) (48 scenes). Tune and test validations both pass.

\begin{table}[t]
\centering
\caption{Holdout validation (recommended \(k=12,b=4\), 48 scenes).}
\label{tab:holdout}
\begin{tabular}{lcccc}
\toprule
Split & FI PSNR \(|\Delta|\) & FI SSIM \(|\Delta|\) & OC PSNR \(|\Delta|\) & OC SSIM \(|\Delta|\) \\
\midrule
Tune & 0.2575 & 0.0005 & 0.1128 & 0.0024 \\
Test & 0.1119 & 0.0039 & 0.1494 & 0.0006 \\
\bottomrule
\end{tabular}
\end{table}

\subsection{Refinement Effect}
Starting from the original 12-scene \(k=6\) subset, metric-aware cluster-preserving refinement reduces objective from 0.2383 to 0.1098, and sharply reduces global gaps (e.g., FI PSNR gap from 0.8441 to 0.0031; OC SSIM gap from 0.0178 to 0.0001).

\section{Discussion}
The results show two practical points. First, risk-aware selection avoids overconfident budget choices by considering uncertainty (LCB), not only average pass rate. Second, enforcing explicit metric-matching constraints keeps subset behavior close to full-data behavior. Split-based holdout evaluation is essential: a subset optimized on one split is not guaranteed to transfer unless tune/test checks are passed.

\section{Limitations}
Our current evaluation uses available OC/FI metric tables as fidelity targets. Full multi-method NeRF ranking fidelity against exhaustive full-dataset runs remains expensive and should be added for final benchmark conclusions. Runtime-aware optimization depends on runtime logs quality; if unavailable, runtime terms are omitted.

\section{Conclusion}
We presented \method, a practical budget-aware subset selection pipeline for NeRF benchmarking at Nutrition5k scale. The system combines balanced candidate generation, risk-aware budget selection, joint \(k\times\)budget search, split-based holdout validation, and metric-aware refinement. In current runs, compact subsets pass strict PSNR/SSIM matching constraints while substantially reducing benchmark size, providing a reproducible pathway toward scalable NeRF benchmarking.

\clearpage

\begin{figure*}[t]
\centering
\includegraphics[width=0.97\linewidth]{figures/fig_pipeline_overview.pdf}
\caption{\method overview.}
\label{fig:pipeline}
\end{figure*}

\begin{figure}[t]
\centering
\includegraphics[width=\linewidth]{figures/fig_budget_pass_lcb.pdf}
\caption{Risk-aware budget behavior.}
\label{fig:budget_curve}
\end{figure}

\begin{figure}[t]
\centering
\includegraphics[width=\linewidth]{figures/fig_holdout_tune_test.pdf}
\caption{Holdout validation summary across FI/OC/ZipNeRF (normalized gaps: \(|\Delta|/\tau\)).}
\label{fig:holdout_curve}
\end{figure}

\begin{figure*}[t]
\centering
\begin{subfigure}[t]{0.49\linewidth}
\centering
\includegraphics[width=\linewidth]{figures/fig_global_gaps.pdf}
\caption{Global FI/OC/ZipNeRF gaps normalized by thresholds (\(|\Delta|/\tau\)).}
\end{subfigure}
\hfill
\begin{subfigure}[t]{0.49\linewidth}
\centering
\includegraphics[width=\linewidth]{figures/fig_refinement_before_after.pdf}
\caption{Refinement effect on FI/OC/ZipNeRF normalized gaps (\(|\Delta|/\tau\)).}
\end{subfigure}
\caption{Additional metric-matching diagnostics from generated report artifacts.}
\label{fig:diagnostics}
\end{figure*}

\begin{thebibliography}{10}\setlength{\itemsep}{-1mm}
\bibitem{mildenhall2020nerf}
B.~Mildenhall, P.~P. Srinivasan, M.~Tancik, J.~T. Barron, R.~Ramamoorthi, and R.~Ng.
\newblock NeRF: Representing scenes as neural radiance fields for view synthesis.
\newblock In \emph{ECCV}, 2020.

\bibitem{barron2021mipnerf}
J.~T. Barron, B.~Mildenhall, D.~Verbin, P.~P. Srinivasan, and P.~Hedman.
\newblock Mip-NeRF: A multiscale representation for anti-aliasing neural radiance fields.
\newblock In \emph{ICCV}, 2021.

\bibitem{muller2022instantngp}
T.~M{\"u}ller, A.~Evans, C.~Schied, and A.~Keller.
\newblock Instant neural graphics primitives with a multiresolution hash encoding.
\newblock \emph{ACM TOG}, 41(4), 2022.

\bibitem{schoenberger2016sfm}
J.~L. Schoenberger and J.-M. Frahm.
\newblock Structure-from-motion revisited.
\newblock In \emph{CVPR}, 2016.

\bibitem{schoenberger2016mvs}
J.~L. Schoenberger, E.~Zheng, M.~Pollefeys, and J.-M. Frahm.
\newblock Pixelwise view selection for unstructured multi-view stereo.
\newblock In \emph{ECCV}, 2016.

\bibitem{gonzalez1985kcenter}
T.~Gonzalez.
\newblock Clustering to minimize the maximum intercluster distance.
\newblock \emph{Theoretical Computer Science}, 38, 1985.

\bibitem{lin2011submodular}
H.~Lin and J.~Bilmes.
\newblock A class of submodular functions for document summarization.
\newblock In \emph{ACL}, 2011.

\bibitem{mirzasoleiman2015lazier}
B.~Mirzasoleiman, A.~Karbasi, R.~Sarkar, and A.~Krause.
\newblock Lazier than lazy greedy.
\newblock In \emph{AAAI}, 2015.
\end{thebibliography}

\end{document}
